\documentclass{article}
\usepackage[T1]{fontenc}
\usepackage{lmodern}
\usepackage{amsmath,amssymb,amsthm,mathtools}
\usepackage[hidelinks]{hyperref}
\usepackage[margin=1in]{geometry}
\newtheorem*{maintheorem}{Theorem}
\newtheorem{lemma}{Lemma}
\DeclareMathOperator{\wt}{wt}
\providecommand{\tightlist}{\setlength{\itemsep}{0pt}\setlength{\parskip}{0pt}}
\title{Analytic distance bounds for every open rectangle in a planar qLDPC family}
\author{}
\date{11 November 2026}
\begin{document}
\maketitle

\begin{maintheorem}[family parameters and distances]
Let \(m,n\) be arbitrary nonnegative integers. The finite binary matrices
\(H_X,H_Z\) defined below on two species of qubits in an
\((m+2)\times(n+2)\) rectangle satisfy
\[
H_XH_Z^{\mathsf T}=0,\qquad
\operatorname{rank}H_X=(m+4)n,\qquad
\operatorname{rank}H_Z=m(n+4).
\]
Their binary CSS logical quotient has dimension eight. Every nontrivial
X logical has Hamming weight at least \(m\), and every nontrivial Z logical
has Hamming weight at least \(n\). Every nontrivial mixed Pauli logical
therefore has weight at least \(\min\{m,n\}\).
\end{maintheorem}

\textbf{Grade: THEOREM.} The complete proof passed three independent hostile
referees, the author and root adversarial rereadings, the corresponding Lean
formalization, and its guarded axiom audits. No numerical certificate,
exhaustive support search, or program-generated case list is used.

The proof has four parts: define the finite matrices and prove their ranks;
prove one-sided exactness by eliminating boundary rows or columns;
prove a uniform polynomial strip inequality; and combine these arguments
with the physical weight identities. All steps used in the main result
are included in this document.
\section{Finite construction and ranks}
\subsection{Definitions}\label{definitions}

The field \(F=\mathbb F_2\) consists of \(0,1\), with addition and
multiplication modulo two. In particular \(a+a=0\) for every \(a\in F\).
Every vector space in this document is over \(F\). A finite vector's
coordinates form a finite index set; its dot product with another vector
is the sum of their coordinatewise products in \(F\).

A finite Laurent polynomial in \(x,y\) is an expression
\(p=\sum  p_{ij}x^iy^j\) with integer exponents and only finitely many
nonzero coefficients. Equality and addition are coefficientwise.
Multiplication is defined by convolution: the coefficient at \((r,s)\)
of \(pq\) is \(\sum _{i+k=r,j+l=s} p_{ij}q_{kl}\). All these sums are
finite. Ordinary polynomials have only nonnegative exponents. A nonzero
coefficient is one. The support is the set of its nonzero coefficient
indices. Write \(\bar p(x,y)=p(x^{-1},y^{-1})\).

The board is

\[
 Q=\{0,1,\ldots,m+1\}\times\{0,1,\ldots,n+1\}.
\]

There are two species \(A,B\) of coordinates at every board cell. Thus
the vector space \(V\) has dimension \(2(m+2)(n+2)\). Write a vector as
a pair \((u,v)\) of polynomials supported on \(Q\), with the first
polynomial representing species A and the second species B. Let \(P\)
delete every coefficient outside \(Q\), separately in each component.

Put

\[
 f=x+x^2+y^2,\qquad g=1+x^2y+x^2y^2.
\]

The X-anchor and Z-anchor sets are

\[
\begin{aligned}
 I_X&=\{-2,-1,\ldots,m+1\}\times\{0,1,\ldots,n-1\},\\
 I_Z&=\{2,3,\ldots,m+1\}\times\{0,1,\ldots,n+3\}.
\end{aligned}
\]

A set of consecutive integers with lower endpoint greater than its upper
endpoint is empty. The matrix rows are

\[
 H_X[(a,b),:]=P\bigl(x^ay^b(f,g)\bigr),\qquad
 H_Z[(c,d),:]=P\bigl(x^cy^d(\bar g,\bar f)\bigr).
\]

The row space of a matrix is the space of all linear combinations of its
rows, and its rank is the dimension of that space. Put
\(S_X=\operatorname{rowspan}(H_X)\) and
\(S_Z=\operatorname{rowspan}(H_Z)\). For a subspace \(S\subseteq V\),
its orthogonal subspace is
\(S^\perp=\{w\in V : s\cdot w=0 \text{for every } s\in S\}\). The
quotient \(S_Z^\perp/S_X\), once \(S_X\subseteq S_Z^\perp\) is
established, is the space of cosets \(w+S_X\) for \(w\in S_Z^\perp\).
For these binary CSS matrices we use ``number of encoded qubits'' to
mean the dimension of this quotient. An \(H_Z\)-cycle is a vector
\(w\) satisfying \(H_Zw=0\), equivalently \(w\in S_Z^\perp\).
An X stabilizer is a vector in \(S_X\); a nontrivial X logical is an
\(H_Z\)-cycle outside \(S_X\). An \(H_X\)-cycle, Z stabilizer, and
nontrivial Z logical are defined by exchanging X and Z in these statements.

\subsection{Polynomial facts}\label{polynomial-facts}

\textbf{Lemma 1.} A product of two nonzero univariate polynomials over
\(F\) is nonzero. The same statement holds for finite Laurent
polynomials in one variable. If a nonzero factor multiplies an unknown
polynomial to zero, the unknown is zero.

\textbf{Proof.} Let \(r\) and \(s\) be the largest exponents with
nonzero coefficients in the two nonzero factors. Every exponent in the
first factor is at most \(r\) and every exponent in the second is at
most \(s\). A pair of such exponents adds to \(r+s\) only when they are
\(r,s\): if the first were less than \(r\), the second would have to
exceed \(s\); if the second were less than \(s\), the first would have
to exceed \(r\). Consequently the coefficient of the product at \(r+s\)
is the product of those two nonzero coefficients, hence one. Thus the
product is nonzero. Finite Laurent supports also have largest exponents,
so that proof applies to them. The final assertion follows by
contradiction if the unknown is nonzero. \hfill\(\square\)

\textbf{Lemma 2.} In \(F[T]\), put \(t=T+T^2\) and
\(b=1+T^2t=1+T^3+T^4\). Then

\[
 b^2+t=1+T+T^2+T^6+T^8\ne0,
 \qquad b^2+T^4t=1+T^5+T^8\ne0.
\]

\textbf{Proof.} Expanding \((1+T^3+T^4)^2\) gives \(1+T^6+T^8\): every
cross term occurs twice and vanishes in \(F\). Adding \(t\) gives the
first displayed polynomial. Since \(T^4t=T^5+T^6\), adding it to
\(1+T^6+T^8\) cancels the two \(T^6\) terms and gives the second. Both
have constant coefficient one, so neither is the zero polynomial. Also
\(t\ne 0\), because its coefficient at exponent one is one.
\hfill\(\square\)

\subsection{Commutation on the finite
board}\label{commutation-on-the-finite-board}

\textbf{Lemma 3.} Every X row is orthogonal to every Z row.

\textbf{Proof.} Choose an allowed X anchor \((a,b)\) and Z anchor
\((c,d)\). Before projection to the board, an A-species overlap pairs an
exponent \(s\) of \(f\) with an exponent \(t\) of \(g\) satisfying

\[
 (a,b)+s=(c,d)-t,
 \quad\hbox{equivalently}\quad s+t=(c-a,d-b).
\]

The parity of the A overlap is therefore the coefficient of \(fg\) at
\((c-a,d-b)\). The B overlap is the coefficient of \(gf\) at the same
exponent: its pairs are an exponent of \(g\) and one of \(f\).
Exchanging the two exponents is a bijection between the two finite sets
of pairs, so these coefficients are equal. Their sum is zero in \(F\).

It remains to show that projection to \(Q\) deletes no overlap. Every
exponent of \(f\) and of \(g\) has both coordinates in \(\{0,1,2\}\). A
point in an unprojected X row has second coordinate between \(b\) and
\(b+2\), hence between \(0\) and \(n+1\), because \(0\le b\le n-1\). A
point in an unprojected Z row has first coordinate between \(c-2\) and
\(c\), hence between \(0\) and \(m+1\), because \(2\le c\le m+1\). A
common point belongs to both rows, so it satisfies both coordinate
bounds and belongs to \(Q\). Thus projection preserves both overlaps and
their sum zero. If either anchor set is empty, the universal
orthogonality assertion has no pair of rows to check. These cases
exhaust all \(m,n\ge 0\): the preceding argument applies whenever the
two selected anchors exist. \hfill\(\square\)

Each entry of \(H_XH_Z^{\mathsf T}\) is the dot product of its selected
X and Z rows. Lemma 3 therefore proves \(H_XH_Z^{\mathsf T}=0\).

\subsection{Independence of every X
row}\label{independence-of-every-x-row}

\textbf{Lemma 4.} The X rows are linearly independent.

\textbf{Proof.} If \(n=0\), the X-anchor set is empty. Its only linear
combination is the empty sum, so the empty family is independent. Now
suppose \(n\ge 1\). An arbitrary linear combination of X rows has an
anchor polynomial

\[
 h(x,y)=\sum_{i=-2}^{m+1} h_i(y)x^i,
 \qquad \operatorname{supp}(h_i)\subseteq\{0,\ldots,n-1\}.
\]

Its physical vector is \(P(fh,gh)\). Suppose this vector is zero. No
coefficient of either product has second coordinate outside
\(\{0,\ldots ,n+1\}\), since every second exponent of \(h\) lies between
\(0,n-1\) and every second exponent of \(f,g\) between \(0,2\). Thus for
every retained first coordinate \(i=0,\ldots ,m+1\), vanishing of all
physical coefficients in that slice is exactly vanishing of these two
ordinary polynomials in \(y\):

\[
 h_{i-2}+h_{i-1}+y^2h_i=0,
 \qquad t h_{i-2}+h_i=0,
 \qquad t=y+y^2.                       \tag{R-X}\label{eq:rank-x}
\]

The second equation gives \(h_i=t h_{i-2}\). Substitution in the first
gives \(h_{i-1}=b h_{i-2}\), where \(b=1+y^2t\). Since \(m\ge 0\), the
retained indices include both \(0\) and \(1\). Applying the substituted
equation at those two indices gives \(h_{-1}=b h_{-2}\) and
\(h_0=b h_{-1}=b^2h_{-2}\). The second equation in~\eqref{eq:rank-x} at \(i=0\) also
gives \(h_0=t h_{-2}\). Hence \((b^2+t)h_{-2}=0\). By Lemma 2 the first
factor is nonzero, so Lemma 1 gives \(h_{-2}=0\). It follows that
\(h_{-1}=0\).

We prove \(h_i=0\) for the remaining indices \(i=0,\ldots ,m+1\) by
increasing-index induction using \(h_i=t h_{i-2}\). For \(i=0\), the
right-hand factor is \(h_{-2}=0\); for \(i=1\), it is \(h_{-1}=0\). For
each \(i\ge 2\), the index \(i-2\) belongs to \(\{0,\ldots ,i-1\}\) and
its slice has already vanished. Thus the equation gives zero in every
case. The two initial slices and these remaining slices exhaust
\(\{-2,\ldots ,m+1\}\). All coefficients of \(h\) vanish, which says
that all coefficients in the original row combination vanish.
\hfill\(\square\)

There are \((m+4)n\) X anchors. Independence makes the rows a basis of
their row space, so their rank is \((m+4)n\).

\subsection{Independence of every Z
row}\label{independence-of-every-z-row}

\textbf{Lemma 5.} The Z rows are linearly independent.

\textbf{Proof.} If \(m=0\), there are no Z anchors and the empty family
is independent. Suppose \(m\ge 1\). Reflect physical cells by

\[
 (i,j)\longmapsto(m+1-i,n+1-j)
\]

and exchange species A and B. Reflection maps \(Q\) bijectively to
itself and is its own inverse; the exchange has the same two properties.
Their combination is therefore a coordinate permutation and maps zero
vectors exactly to zero vectors. An exponent \((c,d)-s\) of a Z row is
taken to \((m+1-c,n+1-d)+s\). Exchanging the species consequently maps
that row to

\[
 P\bigl(x^{m+1-c}y^{n+1-d}(f,g)\bigr).
\]

The transformed first anchor coordinate ranges over
\(\{0,\ldots ,m-1\}\). The transformed second anchor coordinate ranges
over \(\{-2,\ldots ,n+1\}\). Both ranges are bijective images of their
original anchor-coordinate ranges.

A vanishing Z row combination therefore becomes \(P(fh,gh)=0\) with

\[
 h(x,y)=\sum_{j=-2}^{n+1}h_j(x)y^j,
 \qquad \operatorname{supp}(h_j)\subseteq\{0,\ldots,m-1\}.
\]

No first-coordinate truncation occurs in \(fh\) or \(gh\), since the
first exponents lie between \(0\) and \((m-1)+2=m+1\). At each retained
second coordinate \(j=0,\ldots ,n+1\), the zero physical slices are
exactly

\[
 h_{j-2}+t h_j=0,
 \qquad h_j+x^2h_{j-1}+x^2h_{j-2}=0,
 \qquad t=x+x^2.                       \tag{R-Z}\label{eq:rank-z}
\]

Substituting the first equation into the second gives
\(x^2h_{j-1}=b h_j\), where \(b=1+x^2t\). Because \(n\ge 0\), indices
\(0\) and \(1\) occur. The substituted equations at these indices yield

\[
 x^2h_{-1}=b h_0,
 \qquad x^2h_0=b h_1,
 \qquad x^4h_{-1}=b^2h_1.
\]

The first equation of~\eqref{eq:rank-z} at \(j=1\) gives \(h_{-1}=t h_1\), so
\((b^2+x^4t)h_1=0\). Lemmas 2 and 1 give \(h_1=0\). Then
\(x^2h_0=b h_1=0\) implies \(h_0=0\), because \(x^2\) is nonzero. The
first equations of~\eqref{eq:rank-z} at \(j=0,1\) give \(h_{-2}=h_{-1}=0\).

For every remaining index \(j=2,\ldots ,n+1\), increasing-index
induction applies to \(t h_j=h_{j-2}\). The smaller slice has vanished;
since \(t\ne 0\) by Lemma 2, Lemma 1 gives \(h_j=0\). The initial
indices \(-2,-1,0,1\) and these remaining indices exhaust the
transformed anchor range. Every anchor coefficient is zero. Since the
coordinate transformation and the anchor reindexing were bijections,
every coefficient of the original Z row combination is zero.
\hfill\(\square\)

There are \(m(n+4)\) Z anchors, so the Z rank is \(m(n+4)\).

\subsection{Dimension of the logical
quotient}\label{dimension-of-the-logical-quotient}

We give the finite-dimensional linear algebra argument explicitly.

\textbf{Lemma 6.} Let \(T:U\to W\) be a linear map between
finite-dimensional vector spaces. Then
\(\dim  U=\dim  \ker  T+\dim  \operatorname{im} T\). If \(B\subseteq U\)
is a subspace, then \(\dim (U/B)=\dim  U-\dim  B\).

\textbf{Proof.} Choose a basis \(e_1,\ldots ,e_r\) of \(\ker  T\) and
extend it to a basis \(e_1,\ldots ,e_r,v_1,\ldots ,v_s\) of \(U\). Such
an extension is constructed by repeatedly adjoining a vector outside the
span of the current list. Each addition increases the number of
independent vectors by one; this process must stop within \(\dim  U\)
steps. At termination the list spans \(U\) and is independent.

The vectors \(Tv_1,\ldots ,Tv_s\) span \(\operatorname{im} T\): expand
any \(u\in U\) in the displayed basis and apply \(T\), which kills every
\(e_i\). They are independent: a linear combination \(\sum  a_iTv_i=0\)
means \(\sum  a_iv_i\in \ker  T\), hence expresses this combination as a
combination of the \(e_i\). Independence of the extended basis forces
every \(a_i=0\). Thus \(\dim  \operatorname{im} T=s\) and
\(\dim  U=r+s\).

For the quotient assertion, choose a basis of \(B\) and extend it to a
basis of \(U\) by the same construction. The cosets of the added vectors
span \(U/B\), because the original basis vectors lie in \(B\) and hence
have zero coset. A linear combination of the added cosets vanishes
exactly when that combination of added vectors belongs to \(B\);
independence of the extended basis forces every coefficient to vanish.
These cosets are a basis of \(U/B\), proving the asserted dimension.
\hfill\(\square\)

\textbf{Lemma 7.} If independent vectors \(s_1,\ldots ,s_r\) lie in
\(F^D\), then the map \(T:F^D\to F^r\) given by
\(Tw=(s_1\cdot w,\ldots ,s_r\cdot w)\) is onto.

\textbf{Proof.} Suppose its image were proper. Choose a basis of the
image and extend it to a basis of \(F^r\). At least one vector is added.
Assign value one to the first added vector and zero to every other
vector of this extended basis. Extending by linearity gives a nonzero
linear functional \(\ell :F^r\to F\) that vanishes on the image. If
\(u_1,\ldots ,u_r\) are the standard coordinate vectors of \(F^r\) and
\(a_i=\ell (u_i)\), every vector \(z\) has expansion \(\sum z_iu_i\), so
\(\ell (z)=\sum a_iz_i\). At least one \(a_i\) is nonzero, because
\(\ell\) is nonzero. For every \(w\), the identity \(\ell (Tw)=0\) gives
\((\sum a_is_i)\cdot w=0\). Testing successively on every standard
coordinate vector of \(F^D\) shows that every coordinate of
\(\sum a_is_i\) is zero. This contradicts the independence of the
\(s_i\). Therefore the image is all of \(F^r\). If \(r=0\), the target
is the zero vector space and the map is onto by its definition, so the
statement covers that case as well. \hfill\(\square\)

By Lemma 3 every X row is orthogonal to every Z row. Bilinearity of the
dot product gives \(S_X\subseteq S_Z^\perp\), so the quotient in the
rank and dimension assertions exists. Apply Lemma 7 to the independent Z
rows. The corresponding map is precisely \(T(w)=H_Zw\); its kernel is
\(S_Z^\perp\), and its image has dimension \(m(n+4)\). Lemma 6 therefore
gives

\[
 \dim S_Z^\perp=2(m+2)(n+2)-m(n+4).
\]

The X rows are independent, so \(\dim  S_X=(m+4)n\). Apply the quotient
part of Lemma 6 inside \(S_Z^\perp\) to obtain

\[
\begin{aligned}
 \dim(S_Z^\perp/S_X)
 &=2(m+2)(n+2)-m(n+4)-(m+4)n\\
 &=(2mn+4m+4n+8)-(mn+4m)-(mn+4n)\\
 &=8.
\end{aligned}
\]

Every subtraction here subtracts the dimension of a subspace or an image
from that of its ambient space, so no negative dimension is introduced.
For \(m=n=0\), the formula gives \(V=F^8\), \(S_X=S_Z=\{0\}\), and
\(S_Z^\perp/S_X=F^8\). For \(m=0,n>0\), only the Z rows are absent;
Lemma 4 and the same quotient argument apply. For \(n=0,m>0\), only the
X rows are absent; Lemma 5 and the same quotient argument apply.
Together with \(m>0,n>0\), these four mutually exclusive cases exhaust
all nonnegative \(m,n\). This proves the rank and dimension assertions.

\section{Constructive exactness at an open boundary}\label{constructive-exactness-at-an-open-boundary}

For the rectangular code defined above, put \(M=m+2\) and \(N=n+2\).
Every \(H_Z\)-cycle whose two consecutive physical rows \(k,k+1\) are
empty is an X stabilizer. Here \(0\le k\le M-2\). Every \(H_X\)-cycle
whose two consecutive physical columns are empty is a Z stabilizer.

A stronger auxiliary statement holds: an \(H_Z\)-cycle supported in the
first \(r\le M-2\) rows is an X stabilizer, as is one supported in the
last \(r\le M-2\) rows. The analogous assertions hold for Z cycles and
columns.

\subsection{X-sector prefix peeling}\label{x-sector-prefix-peeling}

Write a physical vector as row polynomials \((A_i(T),B_i(T))\) for
\(0\le i<M\), with coefficients in \(\mathbb F_2\) and degrees at most
\(N-1\) whenever nonzero. Put \(s=T^2\) and \(t=T+T^2\). The
\(H_Z\)-cycle equations are precisely

\[
 tA_i+A_{i+2}+B_i+B_{i+1}+sB_{i+2}=0\qquad(0\le i\le M-3).
\]

For any allowed Z anchor \((c,d)\), the dot product of its row with
\((A,B)\) is the coefficient \([x^cy^d](gA+fB)\): a shifted term
\((c,d)-(u,v)\) of \(\bar g\) contributes exactly the coefficient
of A at that coordinate, which is the summand paired with \(x^uy^v\)
in the coefficient of \(gA\). The B terms are paired with the terms of
\(f\) by the identical coefficient formula. Values outside the board
are zero in A and B, exactly as removed by the projection in the matrix row.
The cycle equations therefore follow by reading the coefficient at first exponent
\(i+2\) of \(gA+fB\): \(g=1+x^2(T+T^2)\) and \(f=x+x^2+T^2\). All
transverse coefficients of this syndrome are tested, since the Z-anchor
second coordinates range from \(0\) to \(N+1\).

Suppose the vector is supported in rows \(0,\ldots,r-1\), where
\(1\le r\le M-2\). The cycle equation at index \(i=r-1\) is valid:
\(r-1\le M-3\). Rows \(r\) and \(r+1\) vanish, so this equation says

\[
 B_{r-1}=tA_{r-1}.
\]

Set \(h=A_{r-1}\). If \(h\ne0\), the largest exponent in \(th\) is
\(\deg h+2\), by the nonzero leading coefficient of each factor. Because
\(B_{r-1}\) fits in \(N\) coefficients, \(\deg h+2\le N-1\). Hence
\(\deg h\le N-3\), so \(h\) is a linear combination of the permitted
transverse anchor monomials \(1,T,\ldots,T^{N-3}\). If \(N=2\), a
nonzero \(h\) would give degree at least two but \(B_{r-1}\) has degree
at most one, a contradiction; thus \(h=0\). This handles the empty
transverse anchor set explicitly.

Consider the X stabilizer combination

\[
 p=P\bigl(x^{r-3}h(T)(f,g)\bigr).
\]

Its first anchor coordinate \(r-3\) lies between \(-2\) and \(M-1\); its
transverse anchors are allowed by the previous support argument. Its
physical rows are supported among \(r-3,r-2,r-1\); any negative row is
removed by \(P\). At row \(r-1\), its components are exactly \((h,th)\),
because the \(x^2\) coefficients of \(f,g\) are \(1,t\). It has no
physical rows greater than \(r-1\). Thus subtracting \(p\) removes row
\(r-1\) and leaves support in the first \(r-1\) rows. It remains a
cycle, since X stabilizers are \(H_Z\) cycles by commutation. Induction
on \(r\), with the zero vector at \(r=0\) as base case, expresses the
original vector as a sum of X stabilizers.

\subsection{X-sector suffix peeling}\label{x-sector-suffix-peeling}

Suppose a cycle is supported in rows \(a,\ldots,M-1\), where \(2\le a\le M\).
If \(a=M\), the support is empty and the vector is zero, a stabilizer.
For \(2\le a<M\), the cycle equation with index \(a-2\) is valid and both preceding
physical rows vanish. It gives

\[
 A_a=sB_a.
\]

Set \(h=B_a\). If \(h\ne0\), the largest exponent of \(sh\) is
\(\deg h+2\); therefore \(h\) fits in the permitted transverse anchor
interval \(0,\ldots,N-3\). For \(N=2\) it must be zero. The stabilizer
combination \(P(x^ah(T)(f,g))\) has first row \(a\) with pair
\((sh,h)\), and it is supported only in rows \(a,a+1,a+2\). Its first
anchor coordinate \(a\) lies in the allowed range. Subtracting it
removes row \(a\) and leaves a cycle supported in \(a+1,\ldots,M-1\).
Induction on the number \(M-a\) of remaining rows terminates at the zero
vector.

\subsection{Splitting at two empty
rows}\label{splitting-at-two-empty-rows}

Let rows \(k,k+1\) vanish. Split the vector as \(w=w_L+w_R\), where
\(w_L\) retains only rows below \(k\) and \(w_R\) only rows above
\(k+1\). These sets are disjoint and exhaust the possibly nonzero rows,
so the equality holds coordinatewise.

Each Z check has first-coordinate support contained in three consecutive
rows \(c-2,c-1,c\). It cannot meet both retained parts: such an
intersection would require two rows differing by at least three, while
any two rows in that check differ by at most two. Hence each check pairs
with at most one part, and its zero pairing with \(w\) implies its zero
pairing with both parts. Thus both parts are cycles. The prefix has at
most \(k\le M-2\) rows, and the suffix at most \(M-k-2\le M-2\) rows.
The two peeling arguments put both parts in the X row space, which is
closed under addition.

\subsection{Z-sector orientation and
peeling}\label{z-sector-orientation-and-peeling}

Reflect the physical board and exchange the two species, then describe
the result by columns instead of rows. Let \((A_i(T),B_i(T))\) denote
these column polynomials, with \(0\le i<N\) and transverse degree at
most \(M-1\). A reflected Z stabilizer synthesis with anchor polynomial
sequence \((h_i)\) has components

\[
 A_i=h_i+t h_{i+2},\qquad
 B_i=h_{i+2}+s h_{i+1}+s h_i,
 \qquad t=T+T^2,\ s=T^2.
\]

This is obtained directly by exchanging the variables in \(f,g\) after
the reflection: in the column direction their coefficients are
\((t,0,1)\) for \(f\) and \((1,s,s)\) for \(g\). The anchor index is
shifted by two, so \(h_i\) denotes the anchor at directional coordinate
\(i-2\). Reflection and species exchange are coordinate permutations, so
vanishing and weight are preserved. Under this reflection and exchange,
an X row at \((a,b)\) becomes
\(P(x^{M-1-a}y^{N-1-b}(\bar g,\bar f))\). Its new anchors range
bijectively over \(0\le c\le M+1\), \(2\le d\le N-1\).
Its dot product with the reflected vector is
\([x^cy^d](gA+fB)\), by the coefficient pairing just proved.
For a column expansion, \(g=1+T^2y+T^2y^2\) and
\(f=T+T^2+y^2\). The coefficient at directional exponent \(i+2\)
of this syndrome is
\(A_{i+2}+sA_{i+1}+sA_i+tB_{i+2}+B_i\).
All its transverse coefficients are tested by the stated anchor range.
Thus the actual \(H_X\)-cycle equations in these coordinates are

\[
 tB_{i+2}=A_{i+2}+sA_{i+1}+sA_i+B_i
 \qquad(0\le i\le N-3).
\]

A vector supported in zero columns is zero and is a Z stabilizer.
For a cycle supported in the first \(1\le r\le N-2\) columns, use \(i=r-1\),
which lies in the tested range \(0\le r-1\le N-3\).
The last two physical columns are zero, and the equation gives
\(B_{r-1}=sA_{r-1}\). Therefore \(h=A_{r-1}\) has transverse support in
\(0,\ldots,M-3\) (or is zero when \(M=2\)). A Z synthesis packet with
only \(h_{r-1}=h\) nonzero has last column \(r-1\) equal to \((h,sh)\)
and has no columns beyond it. It removes that column and affects only
the preceding two. This proves prefix exactness by induction.

If the permitted suffix is empty, the vector is zero and is a Z stabilizer.
Otherwise its first potentially nonzero column satisfies \(2\le a<N\);
use \(i=a-2\), which lies in \(0,\ldots,N-3\). The preceding two columns vanish and the equation gives
\(A_a=tB_a\). Thus \(h=B_a\) has allowed transverse support. The
synthesis packet with only \(h_{a+2}=h\) nonzero has first physical
column \(a\) equal to \((th,h)\) and no earlier columns. It removes that
column and affects only the next two. Induction on the number \(N-a\) of allowed suffix columns proves suffix exactness.
To split a Z cycle at two empty columns \(k,k+1\), retain separately the
columns below \(k\) and above \(k+1\). An original X row has second-coordinate
support in \(b,b+1,b+2\); after reflection this still lies in an interval
of length two. It cannot meet both retained parts, whose nearest columns
differ by three. Its pairing with each part is therefore zero because
its pairing with their sum is zero. Both parts are cycles. Their lengths
are \(k\le N-2\) and \(N-k-2\le N-2\), so the two proved Z peeling
statements put them in the Z row space. Their sum lies there too.

\setcounter{lemma}{7}
\section{Uniform polynomial strip inequality}\label{sec:main}

\begin{maintheorem}[arbitrary strip length]
  Let \(L\ge3\) be an integer and let \((a_i,b_i)_{i=0}^{L-1}\) be arbitrary polynomial pairs over \(\mathbb{F}_2\). Assume either of these exact polynomial systems:

  \textbf{Forward system F}, in \(\mathbb{F}_2[y]\):
  \begin{equation}\label{eq:forward}
    a_i = (y + y^2) a_{i-2} + b_{i-1} + b_{i-2} + y^2 b_i \quad (2 \le i \le L-1).
  \end{equation}

  \textbf{Backward system B}, in \(\mathbb{F}_2[x]\):
  \begin{equation}\label{eq:backward}
    (x + x^2) b_i = a_i + x^2 a_{i-1} + x^2 a_{i-2} + b_{i-2} \quad (2 \le i \le L-1).
  \end{equation}

  Let \(\wt(p)\) count the nonzero coefficients of \(p\), with \(\wt(0) = 0\), and put
  \begin{equation}
    w_i = \wt(a_i) + \wt(b_i), \quad W = \sum_{i=0}^{L-1} w_i.
  \end{equation}
  Assume adjacent nonemptiness:
  \begin{equation}\label{eq:adjacent}
    w_i + w_{i+1} \ge 1 \quad (0 \le i \le L-2).
  \end{equation}
  Let \(z\) count indices with \(w_i = 0\), and set \(E = W - (L - z)\). For F define
  \(\varepsilon = 1\) exactly when \(w_{L-3} = w_{L-1} = 0\); for B define \(\varepsilon = 1\) exactly
  when \(w_0 = w_2 = 0\). Otherwise set \(\varepsilon = 0\). Then
  \begin{equation}\label{eq:main-bound}
    E \ge z - 1 - \varepsilon, \quad W \ge L-1 - \varepsilon \ge L-2.
  \end{equation}
\end{maintheorem}


All polynomial equalities are exact, with no quotient reduction, truncation, or degree restriction.

\textbf{Proof organization.} Lemmas~\ref{lem:branch}--\ref{lem:four} establish directional
branching and local gap costs. Lemmas~\ref{lem:propagation}--\ref{lem:exit-alternative} establish
short-run propagation and the run-exit alternative. Section~\ref{sec:accounting} partitions the
excess without double counting and proves the main theorem. Every needed local lemma is proved
here.

\section{Polynomial facts and a branch between zeros}\label{sec:polynomial}

For the local calculations only, rename the indeterminate in either system as \(T\), and put
\(t = T + T^2\). This renames a variable, not a recurrence. A slice is empty if \(w_i = 0\), a
singleton if \(w_i = 1\), and heavy if \(w_i \ge 2\). A gap between consecutive zero indices
\(p < q\) has length \(q - p\) and contains \(q - p - 1\) interior slices. A singleton pair has
exactly one nonzero block, which is a monomial. In particular \(u + T^2 v\) and \(T^2 u + v\) each
have weight one whenever \((u,v)\) is a singleton.

If \(\rho(p)\) counts maximal consecutive runs of exponents in the support of \(p\), with
\(\rho(0) = 0\), then
\begin{equation}\label{eq:runs}
  \wt(tp) = \wt((1 + T)p) = 2\rho(p).
\end{equation}
Indeed a run from exponent \(c\) to \(d\) contributes exactly \(T^c + T^{d+1}\) to \((1 + T)p\).
Endpoints from distinct runs are distinct. A monomial shift does not change weight. Thus
\(\wt(tp) \ge 2\) for \(p \ne 0\). We also use \(\wt(u + v) \le \wt(u) + \wt(v)\) and the fact
that \(p(1)\) is the parity of \(\wt(p)\). Both follow coefficient by coefficient over
\(\mathbb{F}_2\). For two nonzero polynomials with largest exponents \(r, s\), only the pair
\((r,s)\) contributes to the coefficient at \(r + s\) in their product, and that coefficient is 1.
Thus the product is nonzero: the polynomial ring is a domain, and nonzero factors can be
cancelled. If a polynomial vanishes at zero, its constant coefficient is zero; shifting every
remaining exponent down by one expresses it as \(T\) times another polynomial.

\begin{lemma}[directional branch]\label{lem:branch}
  Suppose \(1 \le r \le L-2\) and slices \(r - 1, r + 1\) are empty. In system F, \(b_r = 0\); if
  \(r \le L-3\) and slice \(r\) is nonempty, then \(w_{r+2} \ge 2\). In system B, \(a_r = 0\); if
  \(r \ge 2\) and slice \(r\) is nonempty, then \(w_{r-2} \ge 2\).
\end{lemma}

\begin{proof}[Forward proof]
  Applying equation~\eqref{eq:forward} at \(r + 1\) gives \(0 = b_r\). At \(r + 2\), when this
  index is at most \(L-1\), it gives
  \begin{equation}
    a_{r+2} + T^2 b_{r+2} = t a_r.
  \end{equation}
  The sum of the two physical block weights is at least the weight of their sum, hence at least
  \(2\rho(a_r) \ge 2\) when \(a_r \ne 0\).
\end{proof}

\begin{proof}[Backward proof]
  Applying equation~\eqref{eq:backward} at \(r + 1\) gives \(0 = T^2 a_r\), so \(a_r = 0\).
  Applying equation~\eqref{eq:backward} at \(r\), when \(r \ge 2\), gives
  \begin{equation}
    T^2 a_{r-2} + b_{r-2} = t b_r.
  \end{equation}
  The same coefficient-weight inequality gives \(w_{r-2} \ge 2\rho(b_r) \ge 2\). This proves both
  assertions.
\end{proof}

There is exactly one possible exceptional middle index per direction: \(r = L-2\) in F would
require equation \(L\), and \(r = 1\) in B would require equation 1. Neither equation is among the
hypotheses. No branch cost is asserted at these exceptions.

\section{Local costs inside gaps}\label{sec:gaps}

\begin{lemma}[two occupied slices between zeros]\label{lem:two}
  If \(0 \le k \le L-4\), slices \(k, k + 3\) are empty, and the intermediate slices are not both
  empty, then
  \begin{equation}
    w_{k+1} + w_{k+2} \ge 4.
  \end{equation}
\end{lemma}

\begin{proof}[Forward proof]
  Write the intermediate slices as \((p,q), (u,v)\). Equations at \(k + 2, k + 3\) are
  \begin{equation}
    u = q + T^2 v, \quad tp + q + v = 0.
  \end{equation}
  Evaluation at \(T = 1\) gives \(u(1) = 0\). If \(u \ne 0\), then \(\wt(u) \ge 2\). Also
  \(\wt(p) + \wt(q) + \wt(v) \ge 2\): weight zero forces \(u = 0\), and weight one leaves a single
  nonzero term in \(tp + q + v = 0\). Thus the total is at least four. If \(u = 0\), then
  \(q = T^2 v\) and \(tp = (1 + T^2)v\). Cancelling \(1 + T\) yields \(Tp = (1 + T)v\). Evaluation
  at zero shows that \(v = Th\) for a polynomial \(h\); then \(p = (1 + T)h\) and \(q = T^3 h\).
  Nonemptiness gives \(h \ne 0\), and total weight is \(2\wt(h) + 2\rho(h) \ge 4\).
\end{proof}

\begin{proof}[Backward proof]
  The equations for the same named slices are instead
  \begin{equation}
    tv = u + T^2 p, \quad q = T^2 (u + p).
  \end{equation}
  Evaluation at one gives \(q(1) = 0\). If \(q \ne 0\), its weight is at least two. The remaining
  weight \(\wt(p) + \wt(u) + \wt(v)\) is at least two, since weight zero forces \(q = 0\) and
  weight one cannot satisfy \(tv + u + T^2 p = 0\). If \(q = 0\), then \(u = p\) and
  \(tv = (1 + T^2)p\). Cancelling \(1 + T\) gives \(Tv = (1 + T)p\). Hence \(p = Th\), \(u = Th\),
  \(v = (1 + T)h\) for \(h \ne 0\). The total is again \(2\wt(h) + 2\rho(h) \ge 4\).
\end{proof}

\begin{lemma}[three occupied slices between zeros]\label{lem:three}
  If \(0 \le k \le L-5\), slices \(k, k + 4\) are empty, and all three intermediate slices are
  occupied, then
  \begin{equation}
    w_{k+1} + w_{k+2} + w_{k+3} \ge 4.
  \end{equation}
\end{lemma}

\begin{proof}[Forward proof]
  A failure would make all three pairs \((p,q), (u,v), (h,l)\) singletons. The equations at
  \(k + 2, k + 3, k + 4\) are
  \begin{equation}
    u = q + T^2 v, \quad h = tp + v + q + T^2 l, \quad 0 = tu + l + v.
  \end{equation}
  If \(q = 0\), the first equation forces both \(u, v\) to vanish or both to be nonzero,
  contradicting the second singleton. Thus \(q\) is a monomial and \(p = 0\). If \(u \ne 0\), then
  \(v = 0\) and the last equation gives \(l = tu\) of weight two, contradicting the third
  singleton. Consequently \(u = 0\), \(v\) is a monomial, and \(l = v\), so \(h = 0\). The middle
  equation now gives \(q = (1 + T^2)v\), of weight two, a contradiction.
\end{proof}

\begin{proof}[Backward proof]
  The equations for the three singleton slices are
  \begin{equation}
    tv = u + T^2 p, \quad tl = h + T^2 u + T^2 p + q, \quad v = T^2 (h + u).
  \end{equation}
  If \(v \ne 0\), it is a monomial and \(u = 0\), so the first equation forces \(T^2 p = tv\) of
  weight two, impossible in the first singleton. Hence \(v = 0\) and \(u\) is a monomial. The last
  equation gives \(h = u\), so \(l = 0\). The first gives \(T^2 p = u\), so \(q = 0\). The middle
  equation is then \(0 = u + T^2 u + u = T^2 u\), a contradiction.
\end{proof}

\begin{lemma}[four occupied slices next to a zero]\label{lem:four}
  In system F, an empty slice \(k\), with \(0 \le k \le L-5\), followed by four occupied slices
  satisfies
  \begin{equation}
    \sum_{i=k+1}^{k+4} w_i \ge 5.
  \end{equation}
  In system B, an empty slice \(R\), with \(4 \le R \le L-1\), preceded by four occupied slices
  satisfies
  \begin{equation}
    \sum_{i=R-4}^{R-1} w_i \ge 5.
  \end{equation}
  Neither assertion requires another empty slice at the other end.
\end{lemma}

\begin{proof}[Forward proof]
  If all four slices \((p,q), (u,v), (h,l), (m,n)\) are singletons in increasing index order, the
  equations at \(k + 2, k + 3, k + 4\) are
  \begin{equation}
    u + T^2 v = q, \quad h + T^2 l = tp + v + q, \quad m + T^2 n = tu + l + v.
  \end{equation}
  The first left side is a monomial, so \(q\) is a monomial and \(p = 0\). If the second slice
  were B-only, then \(u = 0, q = T^2 v\), and the second equation would give
  \(h + T^2 l = (1 + T^2)v\) of weight two, a contradiction. Thus the second slice is A-only, with
  \(u = q, v = 0\). The third equation to solve for a slice is now \(h + T^2 l = q\). If the third
  slice is A-only, then \(h = q, l = 0\), and the fourth-slice equation has right side \(tq\), of
  weight two. If it is B-only, then \(h = 0, q = T^2 l\), and that right side is
  \(tq + l = (1 + T^3 + T^4)l\), of weight three. Both contradict the fourth singleton, whose left
  side \(m + T^2 n\) is a monomial.
\end{proof}

\begin{proof}[Backward proof]
  Name singleton slices \((p,q), (u,v), (h,l), (m,n)\) in decreasing order at
  \(R - 1, R - 2, R - 3, R - 4\). The equations at \(R, R - 1, R - 2\) are
  \begin{equation}
    T^2 u + v = T^2 p, \quad T^2 h + l = tq + p + T^2 u, \quad T^2 m + n = tv + u + T^2 h.
  \end{equation}
  The first left side is a monomial, so \(p\) is a monomial and \(q = 0\). If the second slice
  were A-only, then \(u = p, v = 0\), and the second equation would give
  \(T^2 h + l = (1 + T^2)p\) of weight two. Thus the second slice is B-only, with
  \(u = 0, v = T^2 p\). The next equation is \(T^2 h + l = p\). If the third slice is B-only, then
  \(h = 0, l = p\), and the last right side is \(t T^2 p\), of weight two. If it is A-only, then
  \(l = 0, T^2 h = p\), and the last right side is \((1 + T^3 + T^4)p\), of weight three. Neither
  is the monomial \(T^2 m + n\) required by the fourth singleton.
\end{proof}

In a gap of length five between successive zero slices, all four interior slices are occupied.
Lemma~\ref{lem:four} applies from the left endpoint in F and from the right endpoint in B. These
endpoints satisfy \(k \le L-6\) and \(R \ge 5\), respectively. Thus such a gap has interior weight
at least five. All monomials counted in the proofs have distinct exponents. No coefficient may be discarded: the equations hold as ordinary polynomial identities without truncation.

\section{Short runs and exit classification}\label{sec:runs}

Orient the list of zero slices in the forcing direction: increasing index for F and decreasing
index for B. A \emph{short gap} is a gap of length two between successive zero slices. A
\emph{short run} is a maximal sequence of \(m \ge 1\) consecutive short gaps. Its \(m\) middle
slices are nonempty by \eqref{eq:adjacent}. The \emph{exit target} is the next slice after its
last zero in the forcing direction. It either lies inside the following longer gap, lies in the
tail after the last zero, or falls outside the index range.

\begin{lemma}[exact propagation and exceptional exits]\label{lem:propagation}
  Write the active polynomials in the \(m\) middle slices of a short run as
  \(P_1, \dots, P_m\) in forcing order. In F these are A-only slices and in B they are B-only
  slices. Renaming the indeterminate as \(T\) and writing \(t = T + T^2\), one has
  \begin{equation}\label{eq:run-powers}
    P_j = t^{j-1} P_1 \quad (1 \le j \le m), \quad P_1 \ne 0.
  \end{equation}
  At an on-board exit \(e\), respectively,
  \begin{equation}\label{eq:exit-feed}
    a_e + T^2 b_e = t P_m \quad \text{(F)}, \quad
    T^2 a_e + b_e = t P_m \quad \text{(B)}.
  \end{equation}
  Thus every middle after the first, and every on-board exit, has weight at least two. There is
  at most one missing exit: the final middle \(L-2\) in F or the final middle 1 in B. Its occurrence is
  exactly \(\varepsilon = 1\) in the main statement.
\end{lemma}

\begin{proof}
  Lemma~\ref{lem:branch} gives the active block of every middle. In F, at the next middle
  \(r + 2\), the B terms at \(r, r + 1, r + 2\) vanish; hence its A polynomial is \(t\) times that
  at \(r\). At an exit the same equation gives the F identity in \eqref{eq:exit-feed}. In B, the
  equation at a middle \(r\) is \(t b_r = T^2 a_{r-2} + b_{r-2}\) because \(a_r = a_{r-1} = 0\). If
  \(r - 2\) is the next middle, its A polynomial also vanishes, proving the B propagation step;
  otherwise it gives the B exit identity. Induction proves \eqref{eq:run-powers}. The domain
  property and \eqref{eq:runs} show weight at least two at every positive power and every on-board
  exit.

  A run ending at zero \(L-1\) in F has last middle \(L-2\) and would exit at \(L\). Every other F run has an
  on-board exit: in particular zero \(L-2\) has target \(L-1\). In B a run ending at zero 0 has last middle 1
  and would exit at \(-1\); zero 1 still has target 0. Only the final run can lack an exit. Such a
  run exists precisely when the final zero gap is \((L-3,L-1)\) in F or \((2,0)\) in B, as specified
  by \(\varepsilon\).
\end{proof}

\begin{lemma}[cheap exit forces a feed of weight three or four]\label{lem:cheap-exit}
  Suppose three successive exit slices in forcing order have weights \(2, 1, 1\). If their feeding
  middle polynomial is \(P\), then \(\wt(P)\) is three or four. More precisely \(P\) is a monomial
  shift of \(1 + T + T^2\) or \(1 + T + T^2 + T^3\). The following statements specify the exact
  index ranges and equations.
\end{lemma}

\begin{proof}[Forward proof]
  Let the feeding middle be at \(r\), with \(1 \le r \le L-5\), and zeros at \(r - 1, r + 1\). It
  is \((P, 0)\). Write \((A_i, B_i)\) for the slices at \(r + 2 + i\), for \(i = 0, 1, 2\). The
  equations at \(r + 2, r + 3, r + 4\) are
  \begin{equation}\label{eq:cheap-forward}
    \begin{gathered}
      A_0 + T^2 B_0 = tP, \quad A_1 + T^2 B_1 = B_0, \\
      A_2 + T^2 B_2 = t A_0 + B_0 + B_1.
    \end{gathered}
  \end{equation}
  The second equation makes \(B_0\) a monomial; the first slice's weight two then makes \(A_0\) a
  monomial. If the second slice is B-only, the last right side has even coefficient parity:
  \(t A_0\) has even weight and \(B_0, B_1\) are monomials. It cannot be the monomial on the left.
  Thus the second slice is A-only, with \(A_1 = B_0, B_1 = 0\). For \(t A_0 + B_0\) to be a
  monomial, \(B_0\) must cancel exactly one term of \(t A_0\). Hence \(B_0 = T A_0\) or
  \(B_0 = T^2 A_0\). Write \(A_0 = T^a\). The feed equation gives
  \begin{equation}
    P = T^{a-1}(1 + T + T^2) \quad \text{or} \quad
    P = T^{a-1}(1 + T + T^2 + T^3), \quad a \ge 1.
  \end{equation}
  To justify the division, \(a = 0\) would give a nonzero constant on the left of the feed, while
  \(T\) divides its right side. For \(a \ge 1\), cancel \(T(1 + T)\) using
  \(1 + T^3 = (1 + T)(1 + T + T^2)\) and \(1 + T^4 = (1 + T)(1 + T + T^2 + T^3)\). Both displayed
  feed forms have the asserted weight.
\end{proof}

\begin{proof}[Backward proof]
  Let the feeding middle be at \(r\), with \(4 \le r \le L-2\) and zeros at \(r - 1, r + 1\). It is
  \((0, P)\). Write \((A_i, B_i)\) for the slices at \(r - 2 - i\). The equations at
  \(r, r - 1, r - 2\) give
  \begin{equation}\label{eq:cheap-backward}
    \begin{gathered}
      T^2 A_0 + B_0 = tP, \quad T^2 A_1 + B_1 = T^2 A_0, \\
      T^2 A_2 + B_2 = t B_0 + A_0 + T^2 A_1.
    \end{gathered}
  \end{equation}
  Now the second equation makes \(A_0\) a monomial and weight two of the first slice makes \(B_0\)
  a monomial. If the second slice is A-only, then \(A_1 = A_0\) and the last right side is
  \(t B_0 + (1 + T^2)A_0\), of even coefficient parity. It cannot be a monomial. Thus the second
  slice is B-only, and \(t B_0 + A_0\) must be a monomial. Necessarily \(A_0 = T B_0\) or
  \(A_0 = T^2 B_0\). Writing \(B_0 = T^b\), the feed yields
  \begin{equation}
    P = T^{b-1}(1 + T + T^2) \quad \text{or} \quad
    P = T^{b-1}(1 + T + T^2 + T^3), \quad b \ge 1.
  \end{equation}
  If \(b=0\), then \(B_0=1\), and the constant coefficient of
  \(T^2A_0+B_0\) is one, whereas the constant coefficient of \(tP\) is zero.
  Thus \(b\ge1\). The identities
  \(1+T^3=(1+T)(1+T+T^2)\) and
  \(1+T^4=(1+T)(1+T+T^2+T^3)\), followed by cancellation of the nonzero
  factor \(T(1+T)\), give the displayed formulas. This proves the lemma from the separate backward
  equations.
\end{proof}

\begin{lemma}[exit alternative]\label{lem:exit-alternative}
  Suppose a short run feeds a longer zero gap \(D\) of length \(d \ge 4\). Let \(P\) be its final
  middle polynomial and let \(E_D\) be the excess inside \(D\), namely the sum of \(w_i - 1\) over
  its interior. Then
  \begin{equation}\label{eq:exit-alternative}
    E_D \ge 2 \quad \text{or} \quad \wt(P) \ge 3.
  \end{equation}
  For length three, \(E_D \ge 2\) unconditionally.
\end{lemma}

\begin{proof}
  The exit target is the first interior slice and has weight at least two by
  Lemma~\ref{lem:propagation}. If \(E_D \le 1\), that slice has weight exactly two and every other
  interior slice is a singleton. There are at least three interior slices, so
  Lemma~\ref{lem:cheap-exit} applies. Its equations are tested: a gap of length at least four
  requires \(r + 5 \le L-1\) in F or \(r - 5 \ge 0\) in B, stronger than the respective index
  ranges in Lemma~\ref{lem:cheap-exit}. Thus \(\wt(P) \ge 3\). For length three use
  Lemma~\ref{lem:two}.
\end{proof}

\section{Disjoint excess accounting and the main proof}\label{sec:accounting}

If \(z = 0\), all \(L\) slices are occupied, giving \(W \ge L\) and
\(E \ge 0 \ge z - 1 - \varepsilon\). Hence assume \(z \ge 1\). There are \(z - 1\)
successive-zero gaps. All have length at least two by \eqref{eq:adjacent}. Let \(G\) count short
gaps, \(M_g = z - 1 - G\) count longer gaps, and \(R_2\) count maximal short runs. Every interior
slice of a successive-zero gap is occupied, so its contribution to \(E\) is exactly its weight
minus one.

\textbf{Baseline inside runs.} In a short run of \(m\) gaps, reserve one excess unit at every
middle after the first, totalling \(m - 1\). These units exist by Lemma~\ref{lem:propagation}.
Across all short runs this reserves \(G - R_2\) units.

\textbf{Baseline inside longer gaps.} Reserve one excess unit in every longer gap, totalling
\(M_g\). Length three has at least two units by Lemma~\ref{lem:two}. Length four has at least one
by Lemma~\ref{lem:three}. Length at least five has at least one in the first four interior slices
in forcing order by Lemma~\ref{lem:four}. These longer-gap interiors are disjoint from each other
and from all short-run middles. Thus the baseline reserves \(G - R_2 + M_g\) units.

\textbf{An additional unit for each run that has an exit.} If a run is followed by a gap of length
three, that gap has a second excess unit after its reserved baseline. If its following gap has
length at least four, Lemma~\ref{lem:exit-alternative} says either the same second gap unit
exists, or the run's terminal middle has weight at least three. In the latter case that middle
has at least two excess units. At most one was reserved there in the run baseline: one for
\(m \ge 2\), zero for \(m = 1\). At least one further unit therefore remains. Knowing only that
the terminal middle was heavy would not suffice; Lemma~\ref{lem:cheap-exit} supplies the needed
weight at least three.

Distinct short runs have distinct terminal middles and distinct following longer gaps. Run
interiors are disjoint from longer-gap interiors. Thus none of these selected additional units is
a reserved baseline unit or a unit already selected for another run.

If instead the final run exits into an on-board tail, Lemma~\ref{lem:propagation} makes its
target heavy. Its excess unit lies outside all baseline regions and supplies that run's
additional unit. Only a run with a missing exit receives no required additional unit.
Lemma~\ref{lem:propagation} identifies precisely \(\varepsilon\) such runs. Therefore the number
of additional units is at least \(R_2 - \varepsilon\), and
\begin{equation}
  \begin{aligned}
    E &\ge (G - R_2 + M_g) + (R_2 - \varepsilon) \\
      &= G + M_g - \varepsilon = z - 1 - \varepsilon.
  \end{aligned}
\end{equation}
Finally,
\begin{equation}
  W = L - z + E \ge L-1 - \varepsilon \ge L-2.
\end{equation}
This includes \(z = 1\), where \(G = M_g = R_2 = \varepsilon = 0\) and \(W \ge L-1\). The main
theorem is proved. \qed

\textbf{Endpoint distinction.} The only possible discount is the nonexistent target at index \(L\)
in F or \(-1\) in B. An existing target at index \(L-1\) in F or index 0 in B still costs an excess
unit, even when no later zero exists. The result gives \(L-1\) whenever that missing-exit
configuration is absent.


\section{Completion of the physical distance proof}

For a binary vector, define Hamming weight to be the number of its nonzero
coordinates. The support of a polynomial is its set of nonzero coefficient
indices, and its coefficient weight is the size of that support. At a fixed
physical row, every coefficient of the two row polynomials is exactly one
physical bit. The row supports of different rows are disjoint. Consequently
summing the two polynomial coefficient weights over all rows gives exactly
the physical Hamming weight. There is no weight inequality or cancellation
hidden in this identity. The analogous column identity holds after reflection,
exchange of the two species, and exchange of the two coordinate directions,
since those three operations are permutations of the physical coordinates.

Let \(w\in S_Z^\perp\setminus S_X\). Its row polynomials satisfy system F:
the full syndrome coefficient equations proved in the peeling section
rearrange in characteristic two to
\[
 A_i=(T+T^2)A_{i-2}+B_{i-1}+B_{i-2}+T^2B_i
 \quad(2\le i\le M-1),\qquad M=m+2.
\]
There are exactly \(M\) physical row pairs. If two consecutive pairs were
both empty, the constructive exactness proof would put \(w\) in \(S_X\),
contrary to its selection. Therefore every adjacent pair of row weights
sums to at least one. For \(M\ge3\), the uniform strip theorem applies with
\(L=M\) and gives physical weight at least \(M-2=m\).
For \(M=2\), one has \(m=0\), and every physical weight is a nonnegative
integer, giving the same inequality. These are exhaustive cases because
\(M=m+2\ge2\).

Now let \(w\in S_X^\perp\setminus S_Z\). The reflected and exchanged column
polynomials satisfy system B, with exactly \(N=n+2\) columns. The separate
Z-sector peeling proof shows that no adjacent column pair is empty.
For \(N\ge3\), apply the backward strip theorem with \(L=N\), obtaining
weight at least \(N-2=n\). For \(N=2\), one has \(n=0\), and nonnegativity
of weight gives the bound. This proves both claimed sector inequalities.

To state the mixed claim precisely without making a phase convention,
a binary Pauli operator is a pair \((x,z)\in V\times V\), with weight the
number of coordinates \(q\) for which \((x_q,z_q)\ne(0,0)\).
Its symplectic pairing with \((x',z')\) is
\(x\cdot z'+z\cdot x'\). The CSS stabilizer space consists of all pairs
\((s_X,s_Z)\) with \(s_X\in S_X\), \(s_Z\in S_Z\). A pair commutes with
this whole space exactly when \(x\in S_Z^\perp\) and \(z\in S_X^\perp\):
test the pairs \((0,s_Z)\) and \((s_X,0)\) for necessity, and add the two
zero pairings for sufficiency. Such a commuting pair is nontrivial exactly
when \(x\notin S_X\) or \(z\notin S_Z\), because membership in the
Cartesian product requires both memberships.

If \(x\notin S_X\), the X-sector inequality gives \(\wt(x)\ge m\).
Every nonzero coordinate of \(x\) belongs to the support of \((x,z)\),
so \(\wt(x,z)\ge\wt(x)\ge m\ge\min\{m,n\}\).
If \(x\in S_X\), nontriviality forces \(z\notin S_Z\), and the
Z-sector inequality gives
\(\wt(x,z)\ge\wt(z)\ge n\ge\min\{m,n\}\).
These two cases exhaust all nontrivial commuting pairs. Scalar phases do
not change the coordinate pair, its commutation tests, or its weight.
This finishes the main theorem.

\section{Formalization and prior work}
\begingroup\sloppy
The finite construction and ranks are formalized in
\texttt{Planar/RectangleFamily.lean}. The rectangular polynomial encoding
and literal check correspondence are in
\texttt{RectangleFamilySyndrome.lean} and
\texttt{RectangleFamilyMatrix.lean} in the same directory.
\texttt{RectangleFamilyPeeling.lean} and
\texttt{RectangleFamilyGap.lean} formalize the constructive X-sector
exactness argument. \texttt{RectangleFamilyZPeeling.lean} and
\texttt{RectangleFamilyZGap.lean} establish the separate Z-sector argument.
\texttt{RectangleFamilyDistance.lean} connects both physical weights to
the uniform strip theorem and proves the mixed Pauli bound.
The strip theorem is formalized in the three
\texttt{StripFamilyRecurrence}, \texttt{StripFamilyRunCounting}, and
\texttt{StripFamilyRunRecurrence} modules. No finite-matrix evaluation is
an assumption of these family results.

The fixed-polynomial rectangular construction and its eight logical qubits
are established prior work: Liang, Eberhardt, and Chen,
\emph{Planar quantum low-density parity-check codes with open boundaries},
\href{https://arxiv.org/html/2504.08887v4}{arXiv:2504.08887v4},
Figure 1 caption and Section II. That work reports code distances obtained
by integer programming. General boundary descriptions through cellular
automata and polynomial quotients appear in Breuckmann, Choe, Eberhardt,
Pereira, and Steffan,
\emph{Logical Operators and Derived Automorphisms of Tile Codes},
\href{https://arxiv.org/html/2511.14589v1}{arXiv:2511.14589v1}.
The present result supplies a direct analytic proof and Lean verification
of the displayed uniform lower bounds. No priority claim for the
construction or the general boundary-exactness principle is made.
These references identify precedents; no proof step above depends on a
result being supplied by a citation.

\endgroup
\end{document}
